\subsection{Independent proof review: arithmetic determinant transport}\label{endpoint-determinant_review--independent-proof-review-arithmetic-determinant-transport}

Date: 13 September 2026. Review status: the mathematical statements AT1--AT22, including AT2a, are accepted on the corrected source pinned below. The first-pass proof-order issue in the explanation of AT6 has been repaired. This review retains the original Tau Base cohomology, arithmetic theta source, full Taylor unit, tensor sum coordinate, and both reference and arithmetic masses.

\subsubsection{1. Sources and scope}\label{endpoint-determinant_review--sources-and-scope}

The complete first-pass root source was read at SHA-256 41c7b45a0a94f0e65d422b213a4bcbddedbf983ab5127da6095e449e001fef98. Its corrected proof paragraph was then read completely. The accepted source is work/tau\_arithmetic\_determinant\_transport\_20260913.tex, 17706 bytes, SHA-256 24cf7ce5be6a43cc656641b516185df10b2e4d7824600ece657120d1c549deae. No displayed formula changed between those editions.

The original dependencies checked directly were:

{\def\LTcaptype{none} % do not increment counter
\begin{longtable}[]{@{}
  >{\raggedright\arraybackslash}p{(\linewidth - 6\tabcolsep) * \real{0.2308}}
  >{\raggedleft\arraybackslash}p{(\linewidth - 6\tabcolsep) * \real{0.3077}}
  >{\raggedright\arraybackslash}p{(\linewidth - 6\tabcolsep) * \real{0.2308}}
  >{\raggedright\arraybackslash}p{(\linewidth - 6\tabcolsep) * \real{0.2308}}@{}}
\toprule\noalign{}
\begin{minipage}[b]{\linewidth}\raggedright
Source relative to output/split\_zero\_rh\_tandem\_2026-09-12
\end{minipage} & \begin{minipage}[b]{\linewidth}\raggedleft
Bytes
\end{minipage} & \begin{minipage}[b]{\linewidth}\raggedright
SHA-256
\end{minipage} & \begin{minipage}[b]{\linewidth}\raggedright
Scope
\end{minipage} \\
\midrule\noalign{}
\endhead
\bottomrule\noalign{}
\endlastfoot
sources/Tau\_Base\_Cohomology\_2026-09-12/NOTE.tex & 28353 & e7a494d79e01b774ad18045b1e244b28738ddf7b85aec31c58724e2fe35c3aba & Lines 225--408 and 672--755: coefficient sheaf, support faces, projective resolution, actual cohomology, tensor signs and top quotient \\
sources/web\_cyclic\_sum\_delivery/Tau\_Cyclic\_Sum\_Control/NOTE.tex & 41804 & 2cdfd5b25631464f9500dc5fd07c4c8f5c927c8877eb1bf076307ec5dcbc2e8c & Lines 108--183 and 300--349: source, full unit, exact cyclic map and theta primitive \\
tex/arithmetic\_input.tex & 14644 & a50b0fb587b644c4ea94dba45b2b423d038f2bfbe7fef188543939ba850e58d2 & Complete A1--A19, already fully read in this audit \\
\end{longtable}
}

The Gamma norms, all-multiplicity quartet theorem, and exact original arithmetic comparison were independently proved in work/deligne\_bound\_gamma\_independent\_review\_20260913.md, SHA-256 57069ebc4b4eab844d18d3bca9b4b7bd26cda138d9497fe3c126ee65bdfb9186. Its complete 72-equation proof remains a dependency of the application AT22. The section and variation proof for the present transport is given completely in RV6--RV11 below.

The exact finite fixtures described at the end exercise these transport formulas. They are calibration measures for the formulas; they do not assert values of arithmetic zeta moments or the existence of a selected zeta packet. No old suite, Lean process, or remote mutation is part of this review.

\subsubsection{2. Original cohomology and the two separately typed injections}\label{endpoint-determinant_review--original-cohomology-and-the-two-separately-typed-injections}

On the four-point chart ENDPOINTSOURCE0BF632DD06BCE2987885496DSLOT0END, the original coefficient diagram is

ENDPOINTSOURCE0BF632DD06BCE2987885496DSLOT1END

Here ENDPOINTSOURCE0BF632DD06BCE2987885496DSLOT2END is the original even Schwartz space with its value and integral zero, ENDPOINTSOURCE0BF632DD06BCE2987885496DSLOT3END is the original strong Schwartz space on the multiplicative line, ENDPOINTSOURCE0BF632DD06BCE2987885496DSLOT4END, and ENDPOINTSOURCE0BF632DD06BCE2987885496DSLOT5END is the original Fourier transform. The diagram category is the declared linear-fibre realization of the Tau Base chart. At each chart point, the projective resolution

ENDPOINTSOURCE0BF632DD06BCE2987885496DSLOT6END

is exact: at ENDPOINTSOURCE0BF632DD06BCE2987885496DSLOT7END and ENDPOINTSOURCE0BF632DD06BCE2987885496DSLOT8END, its last map is identity; at ENDPOINTSOURCE0BF632DD06BCE2987885496DSLOT9END and ENDPOINTSOURCE0BF632DD06BCE2987885496DSLOT10END, its kernel is the displayed anti-diagonal. Applying its Hom complex to RV1 gives

ENDPOINTSOURCE0BF632DD06BCE2987885496DSLOT11END

Its degree-one image is exactly ENDPOINTSOURCE0BF632DD06BCE2987885496DSLOT12END, so its degree-one cohomology is the original ENDPOINTSOURCE0BF632DD06BCE2987885496DSLOT13END. Tensoring the resolution uses the literal differential sign ENDPOINTSOURCE0BF632DD06BCE2987885496DSLOT14END. At top degree the source is ENDPOINTSOURCE0BF632DD06BCE2987885496DSLOT15END, and the relation subspace is the sum of all tensors with ENDPOINTSOURCE0BF632DD06BCE2987885496DSLOT16END in one factor. Tensoring the quotient maps over the field therefore identifies the top cohomology with ENDPOINTSOURCE0BF632DD06BCE2987885496DSLOT17END. This proves the use of AT1 in the specified category, without replacing the pushforward by evaluation at the zero coefficient stalk ENDPOINTSOURCE0BF632DD06BCE2987885496DSLOT18END.

For an actual full packet ENDPOINTSOURCE0BF632DD06BCE2987885496DSLOT19END, retain ENDPOINTSOURCE0BF632DD06BCE2987885496DSLOT20END, ENDPOINTSOURCE0BF632DD06BCE2987885496DSLOT21END, ENDPOINTSOURCE0BF632DD06BCE2987885496DSLOT22END, and the actual ENDPOINTSOURCE0BF632DD06BCE2987885496DSLOT23END with Mellin transform ENDPOINTSOURCE0BF632DD06BCE2987885496DSLOT24END. The polynomial tensor algebra is ENDPOINTSOURCE0BF632DD06BCE2987885496DSLOT25END, its operator is multiplication by ENDPOINTSOURCE0BF632DD06BCE2987885496DSLOT26END, and its cyclic annihilator is the original ENDPOINTSOURCE0BF632DD06BCE2987885496DSLOT27END.

The injective algebra map ENDPOINTSOURCE0BF632DD06BCE2987885496DSLOT28END sends ENDPOINTSOURCE0BF632DD06BCE2987885496DSLOT29END to ENDPOINTSOURCE0BF632DD06BCE2987885496DSLOT30END. Its kernel before quotienting is exactly the defining cyclic annihilator, proving injectivity. Multiplication by the full tensor unit has inverse multiplication by its inverse. Hence the first injection is

ENDPOINTSOURCE0BF632DD06BCE2987885496DSLOT31END

It is a module map; multiplication by a unit is not being assigned the type of a unital algebra map. The original single-packet section ENDPOINTSOURCE0BF632DD06BCE2987885496DSLOT32END has a left inverse induced by its full arithmetic jet observation. Tensoring this left inverse proves that ENDPOINTSOURCE0BF632DD06BCE2987885496DSLOT33END remains injective. The separate cohomology injection is consequently

ENDPOINTSOURCE0BF632DD06BCE2987885496DSLOT34END

The two target spaces in RV4 and RV5 are carried by the explicit arrow ENDPOINTSOURCE0BF632DD06BCE2987885496DSLOT35END; neither is substituted for the other. Applying the original source map ENDPOINTSOURCE0BF632DD06BCE2987885496DSLOT36END gives exactly ENDPOINTSOURCE0BF632DD06BCE2987885496DSLOT37END and ENDPOINTSOURCE0BF632DD06BCE2987885496DSLOT38END, as in AT2 and AT2a.

Every relation ENDPOINTSOURCE0BF632DD06BCE2987885496DSLOT39END has its original theta primitive. Indeed, the cyclic annihilator says it belongs to the tensor ideal ENDPOINTSOURCE0BF632DD06BCE2987885496DSLOT40END. Fixed-order monic division writes it as ENDPOINTSOURCE0BF632DD06BCE2987885496DSLOT41END. In factor ENDPOINTSOURCE0BF632DD06BCE2987885496DSLOT42END, replace ENDPOINTSOURCE0BF632DD06BCE2987885496DSLOT43END by ENDPOINTSOURCE0BF632DD06BCE2987885496DSLOT44END, apply ENDPOINTSOURCE0BF632DD06BCE2987885496DSLOT45END, and assign the primitive sign ENDPOINTSOURCE0BF632DD06BCE2987885496DSLOT46END. The tensor differential supplies a second ENDPOINTSOURCE0BF632DD06BCE2987885496DSLOT47END; their product is ENDPOINTSOURCE0BF632DD06BCE2987885496DSLOT48END, and ENDPOINTSOURCE0BF632DD06BCE2987885496DSLOT49END recovers the desired source relation. This checks that every variation below retains an actual top cochain boundary.

\subsubsection{3. Finite source, zero relation space, and variations}\label{endpoint-determinant_review--finite-source-zero-relation-space-and-variations}

Let ENDPOINTSOURCE0BF632DD06BCE2987885496DSLOT50END. For ENDPOINTSOURCE0BF632DD06BCE2987885496DSLOT51END, use the original exact sequence

ENDPOINTSOURCE0BF632DD06BCE2987885496DSLOT52END

in increasing monomial frames, with the fixed remainder section ENDPOINTSOURCE0BF632DD06BCE2987885496DSLOT53END. Its coefficient maps do not depend on the real interpolation parameter ENDPOINTSOURCE0BF632DD06BCE2987885496DSLOT54END. The densities are precisely ENDPOINTSOURCE0BF632DD06BCE2987885496DSLOT55END, of mass ENDPOINTSOURCE0BF632DD06BCE2987885496DSLOT56END, and ENDPOINTSOURCE0BF632DD06BCE2987885496DSLOT57END, of mass ENDPOINTSOURCE0BF632DD06BCE2987885496DSLOT58END, in ENDPOINTSOURCE0BF632DD06BCE2987885496DSLOT59END. Neither mass is changed.

Each source density is positive almost everywhere and has moments of every finite degree. For the arithmetic density, the zeros of the nonzero entire function ENDPOINTSOURCE0BF632DD06BCE2987885496DSLOT60END on a real line are discrete; at ENDPOINTSOURCE0BF632DD06BCE2987885496DSLOT61END, its convolution preserves almost-everywhere positivity. A nonzero polynomial has only finitely many zeros on this line, so the two source Grams are positive definite. Their convex interpolation ENDPOINTSOURCE0BF632DD06BCE2987885496DSLOT62END remains positive definite on ENDPOINTSOURCE0BF632DD06BCE2987885496DSLOT63END. At a fixed finite degree, its eigenvalues have a positive lower bound on this compact interval. The relation Gram ENDPOINTSOURCE0BF632DD06BCE2987885496DSLOT64END is positive on its actual domain by injectivity of ENDPOINTSOURCE0BF632DD06BCE2987885496DSLOT65END. These facts justify all finite differentiations, including a small open neighborhood of the interval.

The minimum section is ENDPOINTSOURCE0BF632DD06BCE2987885496DSLOT66END. Orthogonality gives ENDPOINTSOURCE0BF632DD06BCE2987885496DSLOT67END, ENDPOINTSOURCE0BF632DD06BCE2987885496DSLOT68END, and ENDPOINTSOURCE0BF632DD06BCE2987885496DSLOT69END. To prove the inverse expression without presupposing it, decompose every vector ENDPOINTSOURCE0BF632DD06BCE2987885496DSLOT70END uniquely as ENDPOINTSOURCE0BF632DD06BCE2987885496DSLOT71END. Then

ENDPOINTSOURCE0BF632DD06BCE2987885496DSLOT72END

The last equality follows from ENDPOINTSOURCE0BF632DD06BCE2987885496DSLOT73END; it proves ENDPOINTSOURCE0BF632DD06BCE2987885496DSLOT74END. This is the noncircular derivation now in the corrected AT6 proof.

The determinant-line sign is also retained. With ENDPOINTSOURCE0BF632DD06BCE2987885496DSLOT75END, the increasing-degree matrix ENDPOINTSOURCE0BF632DD06BCE2987885496DSLOT76END is triangular with determinant one because ENDPOINTSOURCE0BF632DD06BCE2987885496DSLOT77END is monic. Thus ENDPOINTSOURCE0BF632DD06BCE2987885496DSLOT78END has determinant ENDPOINTSOURCE0BF632DD06BCE2987885496DSLOT79END. The actual line map wedges the ENDPOINTSOURCE0BF632DD06BCE2987885496DSLOT80END relation vectors before the ENDPOINTSOURCE0BF632DD06BCE2987885496DSLOT81END quotient lifts. Replacing a lift by a relation leaves this wedge unchanged. Its Gram block Schur complement is ENDPOINTSOURCE0BF632DD06BCE2987885496DSLOT82END, so

ENDPOINTSOURCE0BF632DD06BCE2987885496DSLOT83END

At ENDPOINTSOURCE0BF632DD06BCE2987885496DSLOT84END, the domain ENDPOINTSOURCE0BF632DD06BCE2987885496DSLOT85END is zero, its determinant is one, and its unique endomorphism is invertible as a map from the zero space to itself. Every composite through that space has the zero map of its stated outer domain and codomain. The formula gives ENDPOINTSOURCE0BF632DD06BCE2987885496DSLOT86END, ENDPOINTSOURCE0BF632DD06BCE2987885496DSLOT87END, and ENDPOINTSOURCE0BF632DD06BCE2987885496DSLOT88END; no earlier singular quotient Gram is inverted.

Differentiation of ENDPOINTSOURCE0BF632DD06BCE2987885496DSLOT89END shows that ENDPOINTSOURCE0BF632DD06BCE2987885496DSLOT90END lies in ENDPOINTSOURCE0BF632DD06BCE2987885496DSLOT91END. Differentiation of ENDPOINTSOURCE0BF632DD06BCE2987885496DSLOT92END fixes its coefficient. Thus

ENDPOINTSOURCE0BF632DD06BCE2987885496DSLOT93END

The two terms involving ENDPOINTSOURCE0BF632DD06BCE2987885496DSLOT94END vanish in the first derivative of ENDPOINTSOURCE0BF632DD06BCE2987885496DSLOT95END by orthogonality. In the next derivative they are the two equal displayed negative terms because ENDPOINTSOURCE0BF632DD06BCE2987885496DSLOT96END is Hermitian. These identities prove AT7 and show that ENDPOINTSOURCE0BF632DD06BCE2987885496DSLOT97END is concave as an original Hermitian form.

Taking determinant derivatives in RV8, or directly in ENDPOINTSOURCE0BF632DD06BCE2987885496DSLOT98END, gives

ENDPOINTSOURCE0BF632DD06BCE2987885496DSLOT99END

The second derivative is exactly

ENDPOINTSOURCE0BF632DD06BCE2987885496DSLOT100END

Each square is for the indicated exact isometry of the original source forms. No term is removed in this identity. It proves AT8 and AT9, including the empty first square when ENDPOINTSOURCE0BF632DD06BCE2987885496DSLOT101END.

Applying ENDPOINTSOURCE0BF632DD06BCE2987885496DSLOT102END to the relation-valued ENDPOINTSOURCE0BF632DD06BCE2987885496DSLOT103END gives the primitive constructed after RV5. The jet ENDPOINTSOURCE0BF632DD06BCE2987885496DSLOT104END and cohomology class ENDPOINTSOURCE0BF632DD06BCE2987885496DSLOT105END are constant in ENDPOINTSOURCE0BF632DD06BCE2987885496DSLOT106END. On a fixed nonempty support label ENDPOINTSOURCE0BF632DD06BCE2987885496DSLOT107END, every such linear map lifts as ENDPOINTSOURCE0BF632DD06BCE2987885496DSLOT108END, with ENDPOINTSOURCE0BF632DD06BCE2987885496DSLOT109END. Thus a nonzero relation vector can map to the supported zero while its energy in RV9--RV11 is retained. Its preimage is the whole source kernel on that label, not only the zero vector.

\subsubsection{4. The exact common-source projection trace}\label{endpoint-determinant_review--the-exact-common-source-projection-trace}

Set ENDPOINTSOURCE0BF632DD06BCE2987885496DSLOT110END, and let ENDPOINTSOURCE0BF632DD06BCE2987885496DSLOT111END be the literal coefficient inclusion. In the common metric ENDPOINTSOURCE0BF632DD06BCE2987885496DSLOT112END, define

ENDPOINTSOURCE0BF632DD06BCE2987885496DSLOT113END

The first two are orthogonal projections with images ENDPOINTSOURCE0BF632DD06BCE2987885496DSLOT114END and ENDPOINTSOURCE0BF632DD06BCE2987885496DSLOT115END. Substitution uses ENDPOINTSOURCE0BF632DD06BCE2987885496DSLOT116END to verify idempotence and ENDPOINTSOURCE0BF632DD06BCE2987885496DSLOT117END-self-adjointness. Since the relation image lies in the polynomial image, ENDPOINTSOURCE0BF632DD06BCE2987885496DSLOT118END. Thus ENDPOINTSOURCE0BF632DD06BCE2987885496DSLOT119END is the orthogonal projection onto ENDPOINTSOURCE0BF632DD06BCE2987885496DSLOT120END, with rank and trace ENDPOINTSOURCE0BF632DD06BCE2987885496DSLOT121END.

The cyclic trace computation has no omitted metric factor:

ENDPOINTSOURCE0BF632DD06BCE2987885496DSLOT122END

Combining these equalities with RV10 proves AT12.

For ENDPOINTSOURCE0BF632DD06BCE2987885496DSLOT123END, set

ENDPOINTSOURCE0BF632DD06BCE2987885496DSLOT124END

For nested subspaces ENDPOINTSOURCE0BF632DD06BCE2987885496DSLOT125END of one Hilbert space, their orthogonal projections satisfy ENDPOINTSOURCE0BF632DD06BCE2987885496DSLOT126END; hence ENDPOINTSOURCE0BF632DD06BCE2987885496DSLOT127END is the projection onto ENDPOINTSOURCE0BF632DD06BCE2987885496DSLOT128END. Each difference in RV14 is therefore a positive projection of rank ENDPOINTSOURCE0BF632DD06BCE2987885496DSLOT129END: the polynomial dimensions differ by ENDPOINTSOURCE0BF632DD06BCE2987885496DSLOT130END; the relation dimensions ENDPOINTSOURCE0BF632DD06BCE2987885496DSLOT131END, including zero at ENDPOINTSOURCE0BF632DD06BCE2987885496DSLOT132END, also differ by ENDPOINTSOURCE0BF632DD06BCE2987885496DSLOT133END. Consequently

ENDPOINTSOURCE0BF632DD06BCE2987885496DSLOT134END

Their sums need not be idempotent, their two summands need not have orthogonal images, and no such assertion enters RV15. Subtracting the four exact traces in RV13 gives

ENDPOINTSOURCE0BF632DD06BCE2987885496DSLOT135END

All inverses are smooth and bounded on the parameter interval. Integrating proves AT15 with the exact arithmetic correction ENDPOINTSOURCE0BF632DD06BCE2987885496DSLOT136END. The individual concavity in RV11 gives no sign to this four-term difference; RV16 is its complete signed expression.

\subsubsection{5. Pointwise expression and arithmetic density zeros}\label{endpoint-determinant_review--pointwise-expression-and-arithmetic-density-zeros}

The two orthogonal source images ENDPOINTSOURCE0BF632DD06BCE2987885496DSLOT137END and ENDPOINTSOURCE0BF632DD06BCE2987885496DSLOT138END decompose the whole finite polynomial space. Their projection formulas imply

ENDPOINTSOURCE0BF632DD06BCE2987885496DSLOT139END

For example, multiplication on the right by ENDPOINTSOURCE0BF632DD06BCE2987885496DSLOT140END acts as identity separately on each of those two images, by orthogonality and their own Gram identities; it therefore acts as identity on their direct sum.

For ENDPOINTSOURCE0BF632DD06BCE2987885496DSLOT141END, evaluation of RV17 gives the exact nonnegative function

ENDPOINTSOURCE0BF632DD06BCE2987885496DSLOT142END

At ENDPOINTSOURCE0BF632DD06BCE2987885496DSLOT143END the relation term is absent and defined as zero. The factor ENDPOINTSOURCE0BF632DD06BCE2987885496DSLOT144END follows from the actual evaluation identity ENDPOINTSOURCE0BF632DD06BCE2987885496DSLOT145END; it is not a replacement of the relation polynomial or its powers.

Integration against ENDPOINTSOURCE0BF632DD06BCE2987885496DSLOT146END yields ENDPOINTSOURCE0BF632DD06BCE2987885496DSLOT147END. Integration against ENDPOINTSOURCE0BF632DD06BCE2987885496DSLOT148END, using the original moment matrices, gives

ENDPOINTSOURCE0BF632DD06BCE2987885496DSLOT149END

Subtracting at the four endpoints and integrating ENDPOINTSOURCE0BF632DD06BCE2987885496DSLOT150END is exactly AT17. To justify both steps without any quotient of densities, note that ENDPOINTSOURCE0BF632DD06BCE2987885496DSLOT151END is a fixed finite-degree polynomial in ENDPOINTSOURCE0BF632DD06BCE2987885496DSLOT152END and ENDPOINTSOURCE0BF632DD06BCE2987885496DSLOT153END, whose coefficients are bounded for ENDPOINTSOURCE0BF632DD06BCE2987885496DSLOT154END. Its absolute value is bounded by a fixed constant times ENDPOINTSOURCE0BF632DD06BCE2987885496DSLOT155END. The inequality ENDPOINTSOURCE0BF632DD06BCE2987885496DSLOT156END and the two actual moment bounds prove absolute integrability on ENDPOINTSOURCE0BF632DD06BCE2987885496DSLOT157END. Fubini applies. At a zero of ENDPOINTSOURCE0BF632DD06BCE2987885496DSLOT158END, this expression is still an ordinary polynomial times a finite signed density; no division by ENDPOINTSOURCE0BF632DD06BCE2987885496DSLOT159END appears. This proves the claimed zero-set and unbounded-ratio domains.

\subsubsection{6. Spectral width and the condition number}\label{endpoint-determinant_review--spectral-width-and-the-condition-number}

The operator ENDPOINTSOURCE0BF632DD06BCE2987885496DSLOT160END is positive and self-adjoint for the original ENDPOINTSOURCE0BF632DD06BCE2987885496DSLOT161END metric. Let ENDPOINTSOURCE0BF632DD06BCE2987885496DSLOT162END be its extreme eigenvalues. Factor the interpolated Gram in its stated matrix order:

ENDPOINTSOURCE0BF632DD06BCE2987885496DSLOT163END

No commutation of the two original Gram matrices is required for this factorization. The matrix on the right of ENDPOINTSOURCE0BF632DD06BCE2987885496DSLOT164END is a rational function of the single diagonalizable positive operator ENDPOINTSOURCE0BF632DD06BCE2987885496DSLOT165END. Its eigenvalue on a ENDPOINTSOURCE0BF632DD06BCE2987885496DSLOT166END-eigenvector is

ENDPOINTSOURCE0BF632DD06BCE2987885496DSLOT167END

The denominator is positive throughout the closed interval. Differentiating with respect to ENDPOINTSOURCE0BF632DD06BCE2987885496DSLOT168END gives ENDPOINTSOURCE0BF632DD06BCE2987885496DSLOT169END, so the extrema are attained at the same ENDPOINTSOURCE0BF632DD06BCE2987885496DSLOT170END for every ENDPOINTSOURCE0BF632DD06BCE2987885496DSLOT171END. Repeated eigenvalues cause no exception. The operator ENDPOINTSOURCE0BF632DD06BCE2987885496DSLOT172END is self-adjoint in ENDPOINTSOURCE0BF632DD06BCE2987885496DSLOT173END because ENDPOINTSOURCE0BF632DD06BCE2987885496DSLOT174END is Hermitian.

For any positive self-adjoint ENDPOINTSOURCE0BF632DD06BCE2987885496DSLOT175END on that same Hilbert space, ENDPOINTSOURCE0BF632DD06BCE2987885496DSLOT176END Indeed ENDPOINTSOURCE0BF632DD06BCE2987885496DSLOT177END is nonnegative: by cyclicity it is the trace of the positive matrix ENDPOINTSOURCE0BF632DD06BCE2987885496DSLOT178END. The lower bound is proved the same way. This does not require ENDPOINTSOURCE0BF632DD06BCE2987885496DSLOT179END to commute with ENDPOINTSOURCE0BF632DD06BCE2987885496DSLOT180END. Applying it to RV15 and subtracting gives

ENDPOINTSOURCE0BF632DD06BCE2987885496DSLOT181END

For ENDPOINTSOURCE0BF632DD06BCE2987885496DSLOT182END, ENDPOINTSOURCE0BF632DD06BCE2987885496DSLOT183END is the derivative of ENDPOINTSOURCE0BF632DD06BCE2987885496DSLOT184END. For ENDPOINTSOURCE0BF632DD06BCE2987885496DSLOT185END, both are respectively zero and the constant zero logarithm, so the same integral gives ENDPOINTSOURCE0BF632DD06BCE2987885496DSLOT186END. Hence

ENDPOINTSOURCE0BF632DD06BCE2987885496DSLOT187END

The triangle inequality for the real integral of RV16, followed by RV22--RV23, proves

ENDPOINTSOURCE0BF632DD06BCE2987885496DSLOT188END

This confirms AT18--AT21 with the factor ENDPOINTSOURCE0BF632DD06BCE2987885496DSLOT189END. The source data at ENDPOINTSOURCE0BF632DD06BCE2987885496DSLOT190END are actual finite Gram data. No bound uniform in packet or tensor degree, and no pointwise upper or lower bound for ENDPOINTSOURCE0BF632DD06BCE2987885496DSLOT191END, is used.

\subsubsection{7. Mass and frame covariance}\label{endpoint-determinant_review--mass-and-frame-covariance}

If ENDPOINTSOURCE0BF632DD06BCE2987885496DSLOT192END with its literal positive scalar ENDPOINTSOURCE0BF632DD06BCE2987885496DSLOT193END, every degree form is ENDPOINTSOURCE0BF632DD06BCE2987885496DSLOT194END, and every minimum section remains the same. Its quotient form is ENDPOINTSOURCE0BF632DD06BCE2987885496DSLOT195END, so its determinant is ENDPOINTSOURCE0BF632DD06BCE2987885496DSLOT196END. The two numerator and two denominator factors in ENDPOINTSOURCE0BF632DD06BCE2987885496DSLOT197END cancel exactly. Thus ENDPOINTSOURCE0BF632DD06BCE2987885496DSLOT198END, also implied by ENDPOINTSOURCE0BF632DD06BCE2987885496DSLOT199END. The separate original masses and determinant factors remain those displayed in these formulas.

For a precise source-frame arrow, let ENDPOINTSOURCE0BF632DD06BCE2987885496DSLOT200END carry new common-source coefficients to old coefficients. Then ENDPOINTSOURCE0BF632DD06BCE2987885496DSLOT201END, ENDPOINTSOURCE0BF632DD06BCE2987885496DSLOT202END, ENDPOINTSOURCE0BF632DD06BCE2987885496DSLOT203END, ENDPOINTSOURCE0BF632DD06BCE2987885496DSLOT204END, and every projection transforms as ENDPOINTSOURCE0BF632DD06BCE2987885496DSLOT205END onto the correspondingly transported subspace. Its inclusion matrix is transported as ENDPOINTSOURCE0BF632DD06BCE2987885496DSLOT206END if the internal frame of ENDPOINTSOURCE0BF632DD06BCE2987885496DSLOT207END is unchanged. Products and traces are then identical by conjugacy. This includes nonunitary complex frames and frames mixing the displayed common-source degree coordinates, because the actual degree subspaces are transported with them.

Independently, let a fixed ENDPOINTSOURCE0BF632DD06BCE2987885496DSLOT208END change the common quotient frame. Then ENDPOINTSOURCE0BF632DD06BCE2987885496DSLOT209END, ENDPOINTSOURCE0BF632DD06BCE2987885496DSLOT210END, and ENDPOINTSOURCE0BF632DD06BCE2987885496DSLOT211END, with the same factor for every ENDPOINTSOURCE0BF632DD06BCE2987885496DSLOT212END. That factor cancels from the four endpoint ratios. These explicit frame maps prove the covariance paragraph following AT21; they do not change the original source norms.

\subsubsection{8. Original quartet application and scope}\label{endpoint-determinant_review--original-quartet-application-and-scope}

For the stipulated actual reflection-stable quartet ENDPOINTSOURCE0BF632DD06BCE2987885496DSLOT213END, with ENDPOINTSOURCE0BF632DD06BCE2987885496DSLOT214END, nonzero imaginary displacement, and complete common order ENDPOINTSOURCE0BF632DD06BCE2987885496DSLOT215END, the independent Gamma theorem retains ENDPOINTSOURCE0BF632DD06BCE2987885496DSLOT216END The exact arithmetic/reference comparison RV24 therefore yields

ENDPOINTSOURCE0BF632DD06BCE2987885496DSLOT217END

Since ENDPOINTSOURCE0BF632DD06BCE2987885496DSLOT218END, moving terms and dividing gives ENDPOINTSOURCE0BF632DD06BCE2987885496DSLOT219END

Both are exactly AT22. They require no existence claim for such a packet. If the fixed packet restriction is invoked, its ENDPOINTSOURCE0BF632DD06BCE2987885496DSLOT220END, Taylor unit, theta primitive, and finite jet injection remain the original ones in RV4--RV6. In particular this is an analytic comparison on the original finite Tau Base complex with its action-compatible arithmetic observation, and does not replace it by a freely chosen polynomial.

The result computes a finite signed arithmetic term and bounds its magnitude using an actual relative Gram spectrum. It neither proves an arithmetic sub-threshold estimate nor draws a conclusion that the entire Tau Base or F1 geometry program has failed.

\subsubsection{9. Exact finite fixtures and final review record}\label{endpoint-determinant_review--exact-finite-fixtures-and-final-review-record}

The independent fixture source and its exact JSON receipt are work/tau\_arithmetic\_determinant\_transport\_exact\_fixtures\_20260913.py and the same basename with .json. The calibration retains the literal original Gamma coordinate ENDPOINTSOURCE0BF632DD06BCE2987885496DSLOT221END and mass ENDPOINTSOURCE0BF632DD06BCE2987885496DSLOT222END, takes the test relation ENDPOINTSOURCE0BF632DD06BCE2987885496DSLOT223END, and compares that Gamma density with ENDPOINTSOURCE0BF632DD06BCE2987885496DSLOT224END. This latter multiplier is positive since it is ENDPOINTSOURCE0BF632DD06BCE2987885496DSLOT225END, and it is unbounded. Its mass is ENDPOINTSOURCE0BF632DD06BCE2987885496DSLOT226END. The test relation is explicitly a calibration polynomial, not a zero-packet assertion.

The fixture tests the first empty relation domain and all four endpoints ENDPOINTSOURCE0BF632DD06BCE2987885496DSLOT227END on cutoff ENDPOINTSOURCE0BF632DD06BCE2987885496DSLOT228END. It checks the exact minimum section, relation orthogonality, full inverse splitting, Schur determinants, common-source traces, both positive-sum traces, and the signed four-endpoint identity at ENDPOINTSOURCE0BF632DD06BCE2987885496DSLOT229END. It explicitly verifies that each sum ENDPOINTSOURCE0BF632DD06BCE2987885496DSLOT230END is not an idempotent and that neither commutes with ENDPOINTSOURCE0BF632DD06BCE2987885496DSLOT231END in this calibration, exercising two assumptions deliberately absent from the proof.

It also differentiates the exact inverse-kernel construction at the first nonzero relation space and checks both matrix variations and both nonzero logarithmic curvature losses. Omitting either loss is detected by a nonzero exact residual. A second fixture rescales the source mass by ENDPOINTSOURCE0BF632DD06BCE2987885496DSLOT232END, checks every quotient scales by that literal scalar and the four-volume correction is zero, and verifies that the relative source operator is ENDPOINTSOURCE0BF632DD06BCE2987885496DSLOT233END. A complex nonunitary frame change additionally verifies the exact similarity of ENDPOINTSOURCE0BF632DD06BCE2987885496DSLOT234END.

One initial harness comparison encountered an unevaluated symbolic zero written ENDPOINTSOURCE0BF632DD06BCE2987885496DSLOT235END. Its exact expression evaluator was repaired; no mathematical identity, hypothesis, or expected value was changed. The actual completed run and source hashes are recorded in the attached receipt. These finite checks supplement the written proof RV1--RV26; they do not prove its analytic moment facts by sampling.

The completed run passed 98 exact checks. Its fixture source SHA-256 is 8f29a721b8409708e3cc6e3b469cf35429743b3e19150fba467636d98085a497, and its JSON output SHA-256 is 367bb7b3c59e092e6b9ea9e444739d83bcf69701e0dd45545dc595f4bbfd67b6. It gives the exact nonzero correction

ENDPOINTSOURCE0BF632DD06BCE2987885496DSLOT236END

At ENDPOINTSOURCE0BF632DD06BCE2987885496DSLOT237END, the positive relation contribution and the positive metric contribution whose negatives occur in RV11 are, respectively,

ENDPOINTSOURCE0BF632DD06BCE2987885496DSLOT238END

Thus this execution checks the two distinct curvature terms with nonzero exact values. No rounding or interval approximation enters these fixture results.

The AT6 proof-order issue was resolved in the accepted source. No further mathematical correction was found. The complete written calculation in RV1--RV28 and the exact fixture receipt support the mathematical acceptance recorded here.
